\section{Polynomial presentation and stabilization}\label{sec:construction}

We construct the counterexample through a hypersurface presentation of a
four-variable polynomial ring.  This
presentation makes stabilization explicit and will also supply the filtration
used to obstruct coordinateness.
It adapts the construction of \cite[Section~2]{Cancellation2026}, with
quadratic terms that permit the polynomial change of coordinates below.

Let
\begin{equation}\label{con:presentation}
 \begin{gathered}
 P=k[p,s,u,F,J],\qquad x=s^2-u^2+pF,\\
 H=x^2F-(1+2sx)J-pJ^2-u,\qquad A=P/(H).
 \end{gathered}
\end{equation}
The five generators of $P$ are independent.  We use the same letters for their
images in $A$; the distinguished element is $p\in A$.

\subsection{A polynomial change of coordinates}

\begin{proposition}\label{prop:polynomial}
There are polynomial coordinates $p',s',F',J$ on $A$ such that
\begin{equation}\label{con:candidate}
 x=(s')^2-J^2+p'F',\qquad
 p=p'-2x\bigl(x(s'+J)+p'J\bigr).
\end{equation}
Thus, under the identification $(p',s',F',J)=(x_1,x_2,x_3,x_4)$,
the element $p$ is the polynomial $f$ in Theorem~\ref{thm:main}.
\end{proposition}

\begin{proof}
The relation $H$ becomes linear after a change of coordinates preserving $x$.
To exhibit that change, write
\[
 M=\begin{pmatrix}F&s-u\\s+u&-p\end{pmatrix},\qquad
 e=\binom{x}{-J},\qquad n=\binom{J}{x}.
\]
Then $x=-\det M$, $e^{\mathsf t}n=0$, and
\[
 H=e^{\mathsf t}Me-J-u.
\]
Keep $J$ fixed and define $F',s',u',p'$ by
\begin{equation}\label{con:matrix-change}
 M'=(I_2-2ne^{\mathsf t})M
   =\begin{pmatrix}F'&s'-u'\\s'+u'&-p'\end{pmatrix}.
\end{equation}
Since $(ne^{\mathsf t})^2=0$ and $\det(I_2-2ne^{\mathsf t})=1$,
we have $-\det M'=x$.  Consequently $x$ and $J$, and hence $e$ and $n$,
are unchanged by the substitution.  Its inverse is the polynomial substitution
\[
 M=(I_2+2ne^{\mathsf t})M',\qquad x=-\det M',
\]
so $p',s',u',F',J$ form a coordinate system on $P$.

Taking half the difference of the off-diagonal entries in
Equation~\eqref{con:matrix-change} gives
\[
 u'=u-e^{\mathsf t}Me,\qquad H=-J-u'.
\]
Eliminating $u'$ therefore identifies $A$ with $k[p',s',F',J]$.
In this quotient $u'=-J$, so the determinant identity gives the first formula
in Equation~\eqref{con:candidate}.  The lower-right entry of the inverse
matrix identity gives
\[
 p=p'-2x\bigl(x(s'-u')+p'J\bigr),
\]
which becomes the second formula after substituting $u'=-J$.
\end{proof}

The expression for $p$ in Equation~\eqref{con:candidate} takes the values $0$ and $1$ at
$(p',s',F',J)=(0,0,0,0)$ and $(1,0,0,0)$, respectively.  In particular,
$p$ is a nonzero nonunit of $A$.

\subsection{One-variable stabilization}

We next change the defining equation after adjoining one variable, while
keeping $p$ fixed.  For this purpose set
\begin{equation}\label{con:xyz}
 y=s+x(x+u^2),\qquad z=sx+pJ.
\end{equation}
Expansion gives the identity
\begin{equation}\label{con:quadric-identity}
 xy-z(z+1)=p(H+u).
\end{equation}
Moreover,
\begin{equation}\label{con:localized-coordinates}
 P[p^{-1}]=k[p,p^{-1}][x,y,z,u],
\end{equation}
with inverse coordinate formulas
\begin{equation}\label{con:localized-inverse}
 s=y-x(x+u^2),\qquad
 F=\frac{x-s^2+u^2}{p},\qquad
 J=\frac{z-sx}{p}.
\end{equation}

\begin{proposition}\label{prop:stable}
For an independent variable $w$, the ring $A[w]$ is a polynomial ring in
five variables with $p$ as one of its coordinates.
\end{proposition}

\begin{proof}
A derivation $\delta$ of a $k$-algebra is \emph{locally nilpotent} if, for
every element $g$, some iterate $\delta^r(g)$ is zero.  We use such a derivation
to replace the equation $H=0$ by $H+pw=0$.

In the coordinates of Equation~\eqref{con:localized-coordinates}, take
\[
 \Delta=-p\frac{\partial}{\partial u},
 \qquad \Delta(p)=\Delta(x)=\Delta(y)=\Delta(z)=0.
\]
This locally nilpotent derivation of $P[p^{-1}]$ preserves $P$, because
Equation~\eqref{con:localized-inverse} yields
\begin{equation}\label{con:delta}
 \begin{gathered}
 \Delta(p)=0,\qquad \Delta(u)=-p,\qquad \Delta(s)=2pxu,\\
 \Delta(F)=-4sxu-2u,\qquad \Delta(J)=-2x^2u.
 \end{gathered}
\end{equation}
Its restriction to $P$ is therefore locally nilpotent as well.  Applying
$\Delta$ to Equation~\eqref{con:quadric-identity} gives $\Delta(H)=p$.
Extend $\Delta$ to $P[w]$ by $\Delta(w)=0$.  The exponential
\[
 \exp(w\Delta)(g)=\sum_{r\geq0}\frac{w^r\Delta^r(g)}{r!}
\]
is a polynomial automorphism: each sum is finite, and its inverse is
$\exp(-w\Delta)$.  It fixes $p$ and sends $H$ to $H+pw$, inducing a
$k[p]$-algebra isomorphism
\begin{equation}\label{con:stabilized-hypersurface}
 A[w]\simeq P[w]/(H+pw).
\end{equation}

It remains to eliminate one variable from the equation on the right.
Put $x_0=s^2-u^2$ and introduce
\begin{equation}\label{con:linear-coordinates}
 \begin{aligned}
 L&=x_0^2F-(1+2sx_0)J,\\
 N&=(1-2sx_0)F+4s^2J.
 \end{aligned}
\end{equation}
The coefficient matrix of this substitution has determinant
$4s^2x_0^2+(1+2sx_0)(1-2sx_0)=1$; explicitly,
\begin{equation}\label{con:linear-inverse}
 F=4s^2L+(1+2sx_0)N,\qquad
 J=-(1-2sx_0)L+x_0^2N.
\end{equation}
Thus $p,s,u,L,N$ are polynomial coordinates on $P$.  Expanding $x=x_0+pF$
in $H$ gives
\begin{equation}\label{con:H-expansion}
 H=L-u+pQ_0,\qquad Q_0=2x_0F^2+pF^3-2sFJ-J^2.
\end{equation}
Since $Q_0\in P$, the substitution $w'=w+Q_0$ is another polynomial change
of coordinates, and
\[
 P[w]/(H+pw)
 =k[p,s,u,L,N,w']/(L-u+pw')
 \simeq k[p,s,u,N,w'].
\]
Together with Equation~\eqref{con:stabilized-hypersurface}, this proves the
assertion with $p$ fixed throughout.
Composing with the identification in Proposition~\ref{prop:polynomial} and
relabeling the target variables so that $p$ maps to $x_1$ gives an
automorphism of $R[w]$ sending $f$ to $x_1$.
\end{proof}

The two steps explicitly realize the exponential shift of
\cite[Lemma~3.3]{DuttaLahiri2021} and the case of the transfer construction
in \cite[Theorem~7]{EdoVenereau2001} in which the perturbation $pQ_0$ is
divisible by $p$.

\subsection{The fibers}

The same coordinate computations describe every fiber of $p$.  This will
give a geometric consequence once noncoordinateness has been proved.

\begin{proposition}\label{prop:fibers}
For every $\lambda\in k$, the quotient $A/(p-\lambda)$ is a polynomial
ring in three variables over $k$.
\end{proposition}

\begin{proof}
At $\lambda=0$, Equations~\eqref{con:linear-coordinates}--\eqref{con:H-expansion}
give
\[
 A/(p)=k[s,u,L,N]/(L-u)=k[s,u,N].
\]
For $\lambda\ne0$, Equations~\eqref{con:localized-coordinates} and
\eqref{con:quadric-identity} give
\[
 A/(p-\lambda)
 =k[x,y,z,u]/\bigl(xy-z(z+1)-\lambda u\bigr)
 \simeq k[x,y,z],
\]
by elimination of $u$.
\end{proof}

Propositions~\ref{prop:polynomial} and~\ref{prop:stable} establish the
polynomial presentation and the stable-coordinate assertion of
Theorem~\ref{thm:main}.  The remaining task is to show that $p$ cannot be
a coordinate of $A$ itself.
